\documentclass[10pt,a4paper]{amsart}
\usepackage[margin=19mm]{geometry}
\usepackage{amsmath,amssymb,mathtools}
\usepackage[scr=rsfs,cal=euler]{mathalfa}
\usepackage{xfrac}
\usepackage[unicode,hidelinks,bookmarks=false]{hyperref}
\newcommand{\ie}{i.e.\ }
\newcommand{\cf}{cf.\ }
\newcommand{\resp}{respectively\ }
\newcommand{\Subsection}{\subsection}
\providecommand{\nobreakdash}{\nobreak\hbox{-}}
\setlength{\emergencystretch}{2em}
\multlinegap0pt
\usepackage{tikz}
\usepackage{amssymb}
\usepackage[all]{xy}
\usepackage{stackengine}
\usepackage{enumerate}
\usepackage{xfrac}
\newtheorem{proposition}{Proposition}
\title{one generator algebras over perfect fields}
\author{Mohamad Maassarani}
\date{Source: arXiv:2512.20187v1 (CC BY 4.0)}
\begin{document}
\maketitle

\noindent\emph{Source and licence.} one generator algebras over perfect fields. Version arXiv:2512.20187v1 (CC BY 4.0), distributed by arXiv under the Creative Commons Attribution 4.0 International licence, http://creativecommons.org/licenses/by/4.0/, as recorded in the arXiv licence metadata for this exact version and shown on the abstract page https://arxiv.org/abs/2512.20187v1. Full licence notice: \texttt{licenses/arxiv-2512.20187-CC-BY-4.0.txt}.\par\smallskip

\noindent\textit{Editorial scope.} Counted unit: the article's proposition proposition (source0.tex lines 183--187). The article's complete original proof of it is delivered with it (source0.tex lines 188--206, one \texttt{proof} environment of 19 lines). No mathematics is written, repaired or re-derived: the statement and the proof are the article's own lines, in the article's own notation, and no retained line is substituted or re-flowed.  No further source-local context is needed: every cross-reference of every retained line resolves inside the two delivered spans.  Nothing was omitted from any retained span, so the package carries no omission record.  No retained line contains a citation, so the wrapper carries no bibliography and no bibliography entry.  No retained line contains a figure environment, an image-inclusion command or drawing code, so no image is delivered and the package declares no asset dependency.  Editorial formatting only, in addition: an AMS \texttt{amsart} preamble that restates the article's own declarations and macros used by the retained lines, namely the environments \texttt{proposition} on separate counters, together with the standard packages those lines need.  The retained environments print in this document as Proposition 1 in order of appearance, whereas the article prints the same items with the numbering of its own full document.  The complete native source of the article as distributed in the arXiv e-print for this exact version is bundled unchanged as \texttt{sources/191544/source0.tex}.
\begin{proposition}
The map $\Theta : Aut(B) \wr \mathfrak{S}_n \to Aut(B^n)$ given by : 
$$ \Theta((f_1,\dots,f_n),\sigma))=(f_1\times \cdots \times f_n)\circ f_\sigma,$$
for $f_1,\dots , f_n\in Aut(B)$ and $\sigma \in \mathfrak{S}_n$ is an injective group morphism.
\end{proposition}

\begin{proof}
We first prove that $\Theta$ is a morphism. For $(f_1,\dots,f_n)$ and $(g_1,\dots, g_n) $ in $Aut(B)^n $, $\sigma$ and $\sigma'$ in $\mathfrak{S}_n$, we have : 
$$\begin{aligned}
\Theta(\ ((f_1,\dots,f_n), \sigma) ((g_1,\dots,g_n),\sigma ') \ ) =&\Theta(\ ((f_1\circ g_{\sigma^{-1}(1)},\dots,f_n\circ g_{\sigma^{-1}(n)}),\sigma\sigma ') \ )\\
=&(f_1\circ g_{\sigma^{-1}(1)}\times \cdots \times f_n\circ g_{\sigma^{-1}(n)}) \circ f_{\sigma\sigma '},
\end{aligned}$$
and 
$$\begin{aligned}
\Theta(\ ((f_1,\dots,f_n), \sigma) \ ) \circ \Theta(\ ((g_1,\dots,g_n),\sigma ') \ ) =&(f_1\times \cdots \times f_n)\circ f_\sigma \circ (g_1\times \cdots \times g_n)\circ f_{\sigma'}\\
=&(f_1\times \cdots \times f_n)\circ (f_\sigma \circ (g_1\times \cdots \times g_n) f_{\sigma^{-1}})\circ f_{\sigma\sigma'} \\
=&(f_1\times \cdots \times f_n)\circ ( g_{\sigma^{-1}(1)}\times \cdots \times g_{\sigma^{-1}(n)})\circ f_{\sigma\sigma'}\\
=&(f_1\circ g_{\sigma^{-1}(1)}\times \cdots \times f_n\circ g_{\sigma^{-1}(n)}) \circ f_{\sigma\sigma '}.
\end{aligned}$$
Comparing the results of both equations, we obtain that $\Theta$ is a morphism. We now prove that $\Theta $ is injective. Asume that $\Theta(\ ((f_1,\dots,f_n), \sigma) \ ) =\text{id}_{B^n}$. This means that for all $b_1,\dots,b_n \in B$, we have 
$$ (f_1(b_{\sigma^{-1}(1)}),\dots, f_n(b_{\sigma^{-1}(n)}))=(b_1,\dots,b_n).$$ 
If $\sigma ^{-1}(i) \neq i$, we get that for all $b\in B$ $a_i=f_i(b)$. This leads to a contradiction because $f_i$ is a bijection. Therefore, $\sigma=\text{id}_{\{1,\dots, n \}}$ and the above equation reduces to  
$$(f_1(b_1),\dots, f_n(b_n))=(b_1,\dots,b_n), $$
for all $b_1,\dots b_n \in B$. This implies that $f_1=\cdots= f_n= \text{id}_B$. We have shown that if $\Theta(\ ((f_1,\dots,f_n), \sigma) \ ) =\text{id}_{B^n}$, then $((f_1,\dots,f_n),\sigma)=((\text{id}_B,\dots, \text{id}_B),\text{id}_{\{1,\dots,n\}})$, proving that $\Theta$ is injective.
\end{proof}

\end{document}
